% Part: computability
% Chapter: computability-theory
% Section: computable-sets

\documentclass[../../../include/open-logic-section]{subfiles}

\begin{document}

\olfileid{cmp}{thy}{cps}
\olsection{గణనీయ సమితులు}

గణనీయ ప్రమేయాల నుంచి గణనీయ సమితులకు గణనీయత భావనను విస్తరించవచ్చు:

\begin{defn}
  $S$ సహజ సంఖ్యల సమితి అనుకుందాం. దాని లక్షణ ప్రమేయం~$\Char{S}$
  గణనీయమైతే, అప్పుడే మాత్రమే $S$ను \emph{గణనీయ సమితి} అంటాం. మరో
  విధంగా చెప్పాలంటే,
\[
\Char{S}(x)=
\begin{cases}
1 & \text{$x\in S$ అయితే} \\
0 & \text{ఇతర సందర్భాల్లో}
\end{cases}
\]
అనే ప్రమేయం గణనీయమైతే, అప్పుడే మాత్రమే $S$ గణనీయం. అదే విధంగా,
$R(x_0,\dots,x_{k-1})$ సంబంధపు లక్షణ ప్రమేయం గణనీయమైతే, అప్పుడే
మాత్రమే ఆ సంబంధం గణనీయం.

గణనీయ సమితులు, సంబంధాలను \emph{నిర్ణయించదగినవి} అని కూడా అంటారు.
\end{defn}

\begin{explain}
పాక్షిక ప్రమేయాలకు, ప్రమేయాలకు, సమితులకు సంబంధించిన అనేక గణనీయత
భావనలు ఇప్పుడు మనకు ఉన్నాయని గమనించండి. వాటిని కలవరపెట్టకండి!{}
\iftag{TMs}{పాక్షిక ప్రమేయాన్ని గణించే ట్యూరింగ్ యంత్రం, ఆ ప్రమేయం
  నిర్వచితమైన ఇన్‌పుట్ విలువలకు దాని అవుట్‌పుట్‌ను తిరిగి ఇస్తుంది;
  సమితిని గణించే ట్యూరింగ్ యంత్రం, ఇన్‌పుట్ విలువ ఆ సమితిలో ఉందో లేదో
  నిర్ణయించి, $1$ లేదా~$0$ను తిరిగి ఇస్తుంది.}{}
\end{explain}

\end{document}
